\section{Limitations and Next Steps}

The proposed method uses the predictive entropy of the agents, and the coupling between them captured by the total correlation, as a proxy for changes in behaviour, with a particular focus on the emergence of communication within a MAS. Its purpose is to provide a channel-agnostic signal, since the improvised message boards used by the agents in the OpenAI--Hugging Face incident were hosted on resources never intended for communication. However, the same breadth that makes the measure channel-agnostic also makes it sensitive to other sources of variation, such as the normal dynamics of agents while solving a task, response length, prompt formatting and temperature. In addition, statistical dependence does not imply communication, since agents sharing a model, prompt or task may be correlated without exchanging any information; joint entropy is hard to estimate with limited samples; and many APIs expose only top-$k$ log-probabilities, so the entropy can only be approximated. The method also requires access to the models' output distributions, which restricts it to the operator hosting the agents, and our experimental evidence is preliminary, based on a single simple benchmark.

Further research is therefore needed to detect specific behaviours from the entropy of the system more precisely. A particularly promising direction is the design and training of a machine learning model on entropy trajectories that distinguishes communication from other behavioural changes. It is also necessary to study more broadly how tools, new information and other perturbations affect the entropy of a MAS, to establish non-communicating baselines across architectures and models, and to validate the approach in more realistic settings, with the aim of building a more robust detector of unintended behaviours and capabilities in general.
